\section{Marco teórico}

Esta sección presenta, de manera sintética, los fundamentos de cada familia de modelos utilizada en el
trabajo y las métricas empleadas para evaluarlos. Dado el volumen de técnicas comparadas, se prioriza una
descripción breve del mecanismo de cada modelo y de su cita de origen, remitiendo a la sección de Análisis
de Resultados para el detalle de los hiperparámetros efectivamente explorados.

\subsection{Modelos ingenuos y estadísticos clásicos}

Las referencias ingenuas (\emph{naive}, \emph{seasonal naive}, \emph{drift} y \emph{average}) constituyen
el piso mínimo de comparación exigido en cualquier ejercicio de pronóstico \parencite{hyndman2021}. La
referencia estacional ingenua con período $m=24$, en particular, replica el último valor observado en la
misma hora del día anterior y resulta especialmente difícil de superar en series con estacionalidad diaria
fuerte, como las estudiadas en este trabajo.

Los modelos ARIMA y su extensión estacional SARIMA combinan componentes autorregresivos, de promedio móvil
y de diferenciación para representar series no estacionarias con estructura de dependencia temporal
\parencite{box2015}. La condición de estacionariedad que estos modelos presuponen se evalúa, en este
trabajo, mediante los contrastes de Dickey-Fuller aumentado \parencite{dickey1979} y de
Kwiatkowski-Phillips-Schmidt-Shin \parencite{kwiatkowski1992}, mientras que la adecuación del ajuste sobre
los residuos se verifica con el estadístico de Ljung-Box \parencite{ljung1978}. El SARIMA ajustado en el
Trabajo Práctico N.\textsuperscript{o} 1 se reutiliza en este trabajo como modelo fundacional, reajustado
(\emph{refit}) sobre el historial ampliado, con el fin de aislar el efecto de la extensión de datos del
efecto de la técnica de modelado. El suavizado exponencial de Holt-Winters y la descomposición estacional
robusta MSTL --que permite múltiples períodos estacionales simultáneos y admite un modelo autorregresivo
sobre la tendencia extraída-- completan el conjunto de técnicas clásicas \parencite{hyndman2021}, ambas
implementadas sobre la librería Statsmodels \parencite{seabold2010}. La selección automática de orden
mediante AutoARIMA, finalmente, recorre un espacio de modelos ARIMA/SARIMA y selecciona el de menor
criterio de información, evitando la identificación manual basada en autocorrelación.

\subsection{Machine Learning: árboles con boosting y ensambles}

Los métodos de \emph{gradient boosting} sobre árboles de decisión ajustan de manera secuencial modelos
débiles que corrigen los errores residuales de la etapa anterior \parencite{friedman2001}. Este trabajo
compara tres implementaciones de uso extendido en la industria: LightGBM, que introduce muestreo por
gradiente y agrupamiento exclusivo de variables para acelerar el entrenamiento en conjuntos de datos
grandes \parencite{ke2017lightgbm}; XGBoost, que formaliza una función objetivo regularizada y una
estrategia eficiente de búsqueda de particiones \parencite{chen2016xgboost}; y CatBoost, que introduce un
esquema de \emph{ordered boosting} para reducir el sesgo de las variables categóricas y del propio proceso
de boosting \parencite{prokhorenkova2018catboost}. Como referencia de ensamble no secuencial se incluye
Random Forest, que promedia árboles entrenados sobre submuestras aleatorias del conjunto de entrenamiento
y de las variables \parencite{breiman2001randomforest}, y Ridge, una regresión lineal regularizada
utilizada como línea de base simple dentro de la familia de Machine Learning. Todos estos modelos se
alimentan con variables de rezago, estadísticos móviles y variables de calendario, generadas y encadenadas
recursivamente mediante la librería skforecast \parencite{amatrodrigo2024skforecast}. Como estrategia de
combinación adicional se evalúa un \emph{stacking} de los tres mejores modelos de esta familia, siguiendo
el principio de generalización apilada de \textcite{wolpert1992stacking}: un metamodelo aprende a
combinar las predicciones de los modelos base a partir de una validación interna.

\subsection{Deep Learning}

Las redes neuronales recurrentes de tipo memoria de largo-corto plazo (LSTM) incorporan compuertas que
regulan el flujo de información a lo largo de la secuencia, mitigando el desvanecimiento del gradiente que
afecta a las redes recurrentes simples \parencite{hochreiter1997lstm}. N-BEATS propone, en cambio, una
arquitectura totalmente conectada, sin componentes recurrentes ni convolucionales, organizada en bloques
residuales que expanden una base interpretable de tendencia y estacionalidad \parencite{oreshkin2020nbeats}.
N-HiTS extiende esa idea con muestreo jerárquico multitasa y una interpolación multiescala orientada a
abaratar el costo computacional en horizontes largos \parencite{challu2023nhits}. Las redes convolucionales
temporales (TCN) aplican convoluciones causales y dilatadas para capturar dependencias de largo alcance sin
recurrencia \parencite{bai2018tcn}. TiDE combina un codificador denso con conexiones residuales para lograr
un costo de entrenamiento reducido frente a arquitecturas basadas en atención, con desempeño competitivo
\parencite{das2023tide}. Finalmente, el Transformer de fusión temporal (TFT) incorpora mecanismos de
atención multihorizonte y selección de variables interpretable, a costa de un tiempo de entrenamiento
sensiblemente mayor \parencite{lim2021tft}. Las seis arquitecturas se implementaron mediante la librería
Darts \parencite{herzen2022darts}.

\subsection{Modelos híbridos y de descomposición}

Prophet formula la serie como una suma de componentes de tendencia, estacionalidad y efectos de
calendario, ajustados mediante un modelo aditivo generalizado con una interfaz orientada a la producción a
escala \parencite{taylor2018prophet}. NeuralProphet retoma esa misma descomposición interpretable pero la
implementa dentro de una red neuronal entrenada por descenso de gradiente, lo que permite incorporar
autorregresión no lineal sobre los rezagos \parencite{triebe2021neuralprophet}. Como estrategia híbrida
adicional se evalúa una descomposición MSTL seguida de un modelo LightGBM sobre la serie
desestacionalizada, que combina la interpretabilidad de la descomposición clásica con la capacidad no
lineal de un modelo de árboles \parencite{hyndman2021,ke2017lightgbm}.

\subsection{AutoML y modelos de fundación}

AutoGluon TimeSeries automatiza la selección y el ensamblado ponderado de un conjunto amplio de modelos
candidatos, sin intervención manual sobre los hiperparámetros \parencite{shchur2023autogluon}. AutoTS sigue
un enfoque análogo, con generaciones evolutivas de combinaciones de modelos sobre una plantilla acotada de
técnicas rápidas. Ambas plataformas se contrastan además con Chronos, un modelo de fundación de series
temporales preentrenado sobre un corpus masivo y heterogéneo de series, evaluado en este trabajo en
modalidad de inferencia directa (\emph{zero-shot}), sin ajuste alguno sobre las series propias
\parencite{ansari2024chronos}. Los modelos estadísticos clásicos y de AutoARIMA se implementaron mediante
la librería StatsForecast \parencite{garza2022statsforecast}.

\subsection{Métricas de error y esquema de validación}

El error absoluto medio (MAE) y la raíz del error cuadrático medio (RMSE) expresan el error en las
unidades originales de la serie, mientras que el error porcentual absoluto medio (MAPE) lo expresa en
términos relativos. Ninguna de estas tres métricas, sin embargo, permite comparar el desempeño entre series
de escalas distintas ni distinguir un buen pronóstico de uno que simplemente aprovecha una estacionalidad
fuerte. El error absoluto escalado medio (MASE) resuelve ambas limitaciones al normalizar el error absoluto
del modelo por el error absoluto, dentro de la muestra de entrenamiento, de una referencia estacional
ingenua de período $m$ \parencite{hyndman2006}; un MASE menor a 1 indica que el modelo supera a esa
referencia. Este trabajo adopta $m=24$ para ambas series, dado que la estacionalidad semanal resultó
comparativamente débil (sección de Análisis de Resultados), y utiliza el MASE de validación como criterio
único de selección del modelo final.

El protocolo de validación empleado consiste en una única ventana de contraste de 48 horas, inmediatamente
anterior al inicio de la intervención de despliegue, y un conjunto de prueba de 105 horas posterior a su
finalización. Esta elección --una sola ventana en lugar de una validación cruzada expandida de múltiples
pliegues-- prioriza mantener acotado el tiempo total de cómputo frente a un conjunto de aproximadamente
quince familias de modelos, a costa de una estimación del error de generalización más ruidosa; esta
limitación se retoma en las Conclusiones. Para evitar la fuga de información, el ajuste de hiperparámetros
y la parada anticipada de las redes profundas se realizan sobre una ventana de 168 horas anterior al
historial de entrenamiento, distinta de la ventana de validación utilizada para clasificar y seleccionar
los modelos.
